\chapter{Intercambio de Información Segura}

\section{Descripción del escenario}

En este apartado se reproduce un intercambio seguro de información entre dos agentes, \textbf{Ana} y \textbf{Berto}, utilizando un 
esquema híbrido que combina:

\begin{itemize}[label=\textbullet]
    \item \textbf{Cifrado simétrico} (AES-256-CBC): Para cifrar el documento (eficiente)
    \item \textbf{Cifrado asimétrico} (RSA): Para proteger las claves simétricas
    \item \textbf{Firma digital} (RSA + SHA-256): Para garantizar integridad y autenticación
\end{itemize}

\subsection{Esquema del protocolo}

El flujo de comunicación es:

\begin{enumerate}
    \item Ana y Berto generan sus pares de claves RSA e intercambian sus claves públicas.
    \item Ana cifra el documento con AES-256-CBC usando una clave y un IV aleatorios.
    \item Ana cifra la clave+IV con la clave pública de Berto (solo Berto puede descifrar).
    \item Ana firma un resumen SHA-256 del documento original con su clave privada.
    \item Ana envía tres ficheros: \texttt{cifrado.txt}, \texttt{claves.txt} y \texttt{firma.txt}.
    \item Berto descifra las claves con su clave privada.
    \item Berto descifra el documento con las claves obtenidas.
    \item Berto verifica la firma de Ana para confirmar integridad y autenticidad.
\end{enumerate}

\begin{figure}[H]
    \centering
    \begin{tabularx}{\textwidth}{|X|}
        \hline
        \textbf{Ana} $\xrightarrow{\text{cifrado.txt + claves.txt + firma.txt}}$ \textbf{Berto} \\
        \hline
        \small{cifrado.txt = AES-256-CBC(texto.txt, clave, IV)} \\
        \small{claves.txt = BASE64(RSA\_Encrypt(clave||IV, $K_{pub\_Berto}$))} \\
        \small{firma.txt = BASE64(RSA\_Sign(SHA256(texto.txt), $K_{priv\_Ana}$))} \\
        \hline
    \end{tabularx}
    \caption{Esquema del intercambio seguro Ana $\rightarrow$ Berto}
\end{figure}

%% ============================================================================
%% Generación de claves
%% ============================================================================

\section{Generación de claves para Ana y Berto}

Se generan dos pares de claves RSA de 2048 bits, protegidos con contraseña:

\begin{lstlisting}[caption=Generación de claves RSA para Ana y Berto]
# Claves de Ana (contraseña: "anak")
openssl genpkey -algorithm RSA -pkeyopt rsa_keygen_bits:2048 \
    -aes-256-cbc -pass pass:"anak" -out anapriv.pem
openssl pkey -in anapriv.pem -passin pass:"anak" \
    -pubout -out anapub.pem

# Claves de Berto (contraseña: "bertok")
openssl genpkey -algorithm RSA -pkeyopt rsa_keygen_bits:2048 \
    -aes-256-cbc -pass pass:"bertok" -out bertopriv.pem
openssl pkey -in bertopriv.pem -passin pass:"bertok" \
    -pubout -out bertopub.pem
\end{lstlisting}

Ana y Berto intercambian sus claves públicas (\texttt{anapub.pem} y \texttt{bertopub.pem}).

%% ============================================================================
%% Trabajo de Ana
%% ============================================================================

\section{Trabajo de Ana: Cifrado y firma}

\subsection{Paso 1: Cifrado del documento}

Ana cifra el fichero \texttt{texto.txt} con AES-256-CBC usando una clave y un IV generados aleatoriamente, sin sal y con salida en BASE64:

\begin{lstlisting}[caption=Generación de clave/IV y cifrado del documento]
# Generar clave (32 bytes) e IV (16 bytes) aleatorios
ANA_KEY=$(openssl rand -hex 32)
ANA_IV=$(openssl rand -hex 16)

# Cifrar con AES-256-CBC (sin sal, salida BASE64)
openssl enc -aes-256-cbc -in texto.txt -out cifrado.txt \
    -K "$ANA_KEY" -iv "$ANA_IV" -nosalt -a
\end{lstlisting}

\subsection{Paso 2: Protección de las claves simétricas}

Ana concatena la clave y el IV en un fichero de texto hexadecimal y lo cifra con la clave pública de Berto:

\begin{lstlisting}[caption=Cifrado de claves simétricas con RSA]
# Crear fichero con clave+IV concatenados
echo "${ANA_KEY}${ANA_IV}" > claves.hex

# Cifrar con la clave pública de Berto (salida binaria)
openssl pkeyutl -encrypt -in claves.hex \
    -pubin -inkey bertopub.pem -out claves.bin

# Convertir a BASE64
openssl enc -a -in claves.bin -out claves.txt
\end{lstlisting}

\textbf{Formato de \texttt{claves.hex}:} Los primeros 64 caracteres hexadecimales corresponden a la clave AES-256 (32 bytes) y 
los siguientes 32 caracteres al IV (16 bytes).

\subsection{Paso 3: Firma digital}

Ana obtiene el resumen SHA-256 del documento original y lo firma con su clave privada:

\begin{lstlisting}[caption=Firma digital del documento]
# Obtener resumen SHA-256 en binario
openssl dgst -sha256 -binary -out resumen.bin texto.txt

# Firmar con clave privada de Ana
openssl pkeyutl -sign -in resumen.bin \
    -inkey anapriv.pem -passin pass:"anak" \
    -out firma.bin

# Convertir a BASE64
openssl enc -a -in firma.bin -out firma.txt
\end{lstlisting}

\subsection{Envío}

Ana envía a Berto los tres ficheros junto con la meta-información:

\begin{table}[H]
    \centering
    \caption{Ficheros enviados por Ana a Berto}
    \begin{tabularx}{\textwidth}{l|l|p{6cm}}
        \toprule
        \textbf{Fichero} & \textbf{Formato} & \textbf{Contenido} \\
        \midrule
        \texttt{cifrado.txt} & BASE64 & Documento cifrado con AES-256-CBC \\
        \texttt{claves.txt} & BASE64 & Clave+IV cifrados con RSA (clave pública de Berto) \\
        \texttt{firma.txt} & BASE64 & Firma SHA-256 con RSA (clave privada de Ana) \\
        \bottomrule
    \end{tabularx}
\end{table}

\textbf{Meta-información:} Algoritmo AES-256-CBC, concatenación clave (64 hex) + IV (32 hex).

\section{Trabajo de Berto: Descifrado y verificación}

\subsection{Paso 1: Recuperación de las claves simétricas}

Berto convierte \texttt{claves.txt} de BASE64 a binario y descifra con su clave privada:

\begin{lstlisting}[caption=Recuperación de claves por Berto]
# Convertir BASE64 a binario
openssl enc -a -d -in claves.txt -out claves2.bin

# Descifrar con clave privada de Berto
openssl pkeyutl -decrypt -in claves2.bin \
    -inkey bertopriv.pem -passin pass:"bertok" \
    -out claves2.hex
\end{lstlisting}

El fichero \texttt{claves2.hex} debe ser idéntico a \texttt{claves.hex} de Ana.

\subsection{Paso 2: Extracción de clave e IV}

Siguiendo la meta-información proporcionada por Ana:

\begin{lstlisting}[caption=Extracción de clave e IV]
# Leer contenido
CONTENIDO=$(cat claves2.hex | tr -d '\n')

# Clave AES-256: primeros 64 caracteres hex
BERTO_KEY="${CONTENIDO:0:64}"

# IV: siguientes 32 caracteres hex
BERTO_IV="${CONTENIDO:64:32}"
\end{lstlisting}

\subsection{Paso 3: Descifrado del documento}

\begin{lstlisting}[caption=Descifrado del documento por Berto]
openssl enc -aes-256-cbc -d -in cifrado.txt \
    -out mensaje2.txt \
    -K "$BERTO_KEY" -iv "$BERTO_IV" -nosalt -a
\end{lstlisting}

El fichero \texttt{mensaje2.txt} resultante debe ser \textbf{idéntico} al \texttt{texto.txt} original de Ana.

\subsection{Paso 4: Verificación de la firma}

La verificación se realiza mediante \texttt{-verifyrecover}, que permite recuperar el resumen original que Ana firmó:

\begin{lstlisting}[caption=Verificación de la firma digital]
# Convertir firma de BASE64 a binario
openssl enc -a -d -in firma.txt -out firma2.bin

# Recuperar el resumen de Ana (descifrar firma con clave publica)
openssl pkeyutl -verifyrecover -in firma2.bin \
    -pubin -inkey anapub.pem -out resumen_ana.bin

# Calcular resumen SHA-256 del mensaje descifrado
openssl dgst -sha256 -binary -out resumen_berto.bin mensaje2.txt

# Comparar ambos resumenes
diff resumen_ana.bin resumen_berto.bin
\end{lstlisting}

\textbf{Resultado esperado:}

\begin{itemize}[label=\textbullet]
    \item Si \texttt{diff} no produce salida: Los resúmenes son \textbf{idénticos}.
    \item \textbf{Integridad confirmada:} El documento no ha sido modificado durante la transmisión.
    \item \textbf{Autenticidad confirmada:} Solo Ana (con su clave privada) pudo generar esa firma.
\end{itemize}

\section{Análisis de seguridad del protocolo}

\begin{table}[H]
    \centering
    \caption{Propiedades de seguridad del esquema híbrido}
    \begin{tabularx}{\textwidth}{l|l|p{6cm}}
        \toprule
        \textbf{Propiedad} & \textbf{Mecanismo} & \textbf{Justificación} \\
        \midrule
        Confidencialidad & AES-256-CBC + RSA & Solo Berto puede descifrar la clave simétrica \\
        Integridad & SHA-256 & El resumen detecta cualquier modificación \\
        Autenticación & Firma RSA & Solo Ana posee la clave privada para firmar \\
        No repudio & Firma RSA & Ana no puede negar haber firmado el documento \\
        \bottomrule
    \end{tabularx}
\end{table}

\subsection{Limitaciones}

\begin{itemize}[label=\textbullet]
    \item \textbf{Sin certificados:} No hay una autoridad de certificación que garantice la identidad. Ana y Berto confían directamente
     en las claves intercambiadas.
    \item \textbf{Sin protección contra ataques de repetición ni man in the middle:} Un atacante podría reenviar los mismos tres ficheros.
    \item \textbf{Sin forward secrecy:} Si la clave privada de Berto se compromete en el futuro, comunicaciones pasadas pueden ser 
    descifradas.
\end{itemize}
